% Section 8 (part 5): Theorem P(a), parts (ii)-(iv); the uniformising map; reduction modulo exact forms.

\begin{proof}[Proof of (ii)]
Let $A:=\bbQ\lser{s}$, a field, in which $L=s^{-1}$ and $u=s^5$. By \cref{lem:fsp-univ}, $\jmath(c)$ is the normalised saddle series of the same data computed over $A$. As in the proof of (i), $\mathsf S_0$ stands for the formal phase $\mathsf S_0(L+t)-\mathsf S_0(L)$ and the amplitude is $a_0:=A_{\mathrm{St}}/A_{\mathrm{St}}(L,0)$, both with Taylor data in $\AL$, mapped to $A$ by $\jmath$. The letters $\phi$, $J$, $a$, $a_0$ are local to this proof.

\emph{Step 1 (equivariance).} On $A$, $\Dq=-(1-u)\,\dd/\dd u=-\frac{1-s^5}{5s^4}\frac{\dd}{\dd s}$, and $\Dq(s^{-1})=(1-s^5)/(5s^6)=\jmath((L^6-L)/5)$. A ring homomorphism compatible with the derivations on the generators is compatible on the localisation; so $\jmath\circ\Dq=\Dq\circ\jmath$.

\emph{Step 2 (change of variables).} The map $\phi(X):=L/X$, $\phi(X)-L=L\sum_{m\geq1}(1-X)^m$, has $\phi(1)=L$ and $J:=\phi'(1)=-L$, a unit of $A$. By \cref{lem:fsp-cov} over $A$, $\SPn_L[\mathsf S_0;a]=\SPn_{X=1}[\mathsf S_0\circ\phi;(a\circ\phi)\,\phi'/J]$ for every amplitude $a\in A\fps{\sigma-L}\lser z$.

\emph{Step 3 (the phases agree).} In $A\fps{X-1}$, $\mathsf S_0(\phi(X))-\mathsf S_0(L)=-s^4\bigl(\hat\cW(X;u)-\hat\cW(1;u)\bigr)$. Both sides vanish at $X=1$, and their $X$-derivatives are $-(L/X^2)\mathsf S_0'(L/X)$ and $(L/X^2)\log((1-uX^5)/(1-u))$ (\cref{lem:con-form}~(c)); so it suffices that $\mathsf S_0'(L/X)=-\log((1-uX^5)/(1-u))$. Both sides of this vanish at $X=1$, and both $X$-derivatives are images of $5uX^4/(1-uX^5)$, $u=s^5$ (on the left $\mathsf S_0''(L/X)\cdot(-L/X^2)$, with $\mathsf S_0''=-5/(\sigma(\sigma^5-1))$), under two ring homomorphisms $\bbQ[s^{\pm1},X^{\pm1},(1-s^5X^5)^{-1}]\to A\fps{X-1}$: expansion at $X=1$ over $\bbQ(s)$, and expansion of $(1-s^5X^5)^{-1}=\sum_Ps^{5P}X^{5P}$ in $\bbQ[X,X^{-1}]\lser s$ followed by expansion at $X=1$ coefficientwise in $s$, a ring homomorphism $\bbQ[X,X^{-1}]\lser s\to A\fps{X-1}$ which commutes with $\dd/\dd X$. They agree on the generators, hence everywhere; and a series in $X-1$ is determined by its value at $X=1$ and its derivative.

\emph{Step 4 (the amplitudes agree).} In $A\fps{X-1}\fps{z}$, with $z=s^4\zh$, $(A_{\mathrm{St}}\circ\phi)\,\phi'/(J\,A_{\mathrm{St}}(L,0))=\hat\fa/\hat\fa(1;0)$. The half-density: at $\sigma=L/X$ one has $5\sigma/(1-\sigma^5)=5LuX^4/(uX^5-1)$ and $\phi'/J=X^{-2}$, so $\re^{\mathsf S_1(L/X)-\mathsf S_1(L)}\phi'/J$ is the square root with constant term $1$ of $(1-u)/(1-uX^5)$, which is the half-density of $\hat\fa/\hat\fa(1;0)$. The Stirling parts: $z^{2k-1}(5\sigma)^{1-2k}=5^{1-2k}(uX\zh)^{2k-1}$ and
\[
z^{2k-1}\sum_j(\zeta^j-\sigma)^{1-2k}=-(uX\zh)^{2k-1}\sum_j(1-\zeta^jsX)^{1-2k}=-(uX\zh)^{2k-1}\,5\sum_{P\geq0}\binom{2k-2+5P}{5P}u^PX^{5P},
\]
which gives the exponent of $\hat\fa$ (\cref{lem:con-form}~(b), (c)); these identities hold by the two-homomorphism argument of Step~3, applied to the sums over $j$, which lie in $\bbQ[s,X,(1-s^5X^5)^{-1}]$.

\emph{Step 5 (the rescaling).} Write the amplitude of Step 4 as $a(X;z):=\hat\fa(X;z/s^4)/\hat\fa(1;0)$. By \cref{lem:fsp-homog} with the unit $\mu=s^4$ of $A$, $\SPn_{X=1}[-s^4\hat\cW;a(\cdot;z)](z)=\SPn_{X=1}[-\hat\cW;a(\cdot;s^4\zh)](\zh)$ at $\zh=z/s^4$, and $a(\cdot;s^4\zh)=\hat\fa(\cdot;\zh)/\hat\fa(1;0)$. With Steps 2--4 and \cref{lem:fsp-const} this gives $\jmath(c)(z)=\PhiSP(z/s^4;u)$ in $A\fps{z}$. The hatted data lie in $\bbQ\fps{u}$ (\cref{lem:con-form}~(c)), so $\PhiSP$ is computed in $\bbQ\fps{u}\fps{\zh}$ and, by \cref{lem:fsp-univ} applied to $\bbQ\fps{u}\to A$, equals the same functional over $A$.
\end{proof}

\begin{proof}[Proof of (iii)]
Apply $\jmath$ to (i): by Step 1 of the proof of (ii), $\PFfive^{L}(s^{-1}\jmath(c))=0$ in $\bbQ\lser{s}\fps{z}$. The substitution $z=s^4\zh$ is a ring isomorphism $\bbQ\lser{s}\fps{z}\to\bbQ\lser{s}\fps{\zh}$, since $s^4$ is a unit, and by (ii) it maps $s^{-1}\jmath(c)$ to $s^{-1}\PhiSP$. By \eqref{eq:con-chain}, $\PFfive^{L}(s^{-1}\PhiSP)=s^{-1}\mathrm{HPF}(\PhiSP)$, so $\mathrm{HPF}(\PhiSP)=0$ in $\bbQ\lser{s}\fps{\zh}$. HPF has coefficients in $\bbQ[u,\zh]$, and $\PhiSP$ lies in $\bbQ\fps{u}\fps{\zh}$, which embeds in $\bbQ\lser{s}\fps{\zh}$; so the identity holds in $\bbQ\fps{u}\fps{\zh}$.
\end{proof}

\begin{proof}[Proof of (iv)]
By (iii), $\PhiSP$ solves HPF with $\mathrm{ph}^{\mathrm{SP}}_0=1$ and coefficients in $\bbQ\fps{u}$; by \cref{lem:gau-poly} (only its polynomiality clause is used) every $\mathrm{ph}^{\mathrm{SP}}_n$ is a polynomial. Then $\jmath\bigl(L^{4n}\mathrm{ph}^{\mathrm{SP}}_n(L^{-5})\bigr)=s^{-4n}\mathrm{ph}^{\mathrm{SP}}_n(s^5)=\jmath(c_n)$ by (ii), and $\jmath$ is injective.
\end{proof}

Theorem P(b) identifies $\PhiSP$ with $\Phican$. Its proof uniformises the hatted phase by the model phase $\Wv$, reduces the $1$-form modulo exact forms to the five model forms, and compares two functionals on the reduced form, the stationary phase at the critical point and the residue at the pole.

\begin{lem}[The uniformising map]\genusstatus{AI-proposed, unchecked}
\label{lem:con-cfu}
There are unique $Y_k\in\bbQ[T]$ ($k\geq1$) and $V\in1+u\bbQ\fps{u}$ such that $Y:=X\bigl(1+\sum_{k\geq1}u^kY_k(X^5)\bigr)$ satisfies
\[
\hat\cW(X;u)=V/Y+Y^4/4=W_V(Y)\qquad\text{in }\bbQ[X,X^{-1}]\fps{u}.
\]
Moreover $\deg Y_k\leq k$ and $Y_1=(1+T)/18$; $V=\tau^{-5/4}$; the point $X=1$ goes to the critical point $Y_c:=Y(1)=V^{1/5}=\tau^{-1/4}$ of $W_V$, and $Y'(1)\in1+u\bbQ\fps{u}$; the inverse map has the same form, $X=Y\bigl(1+\sum_ku^kX_k(Y^5)\bigr)$ with $X_k\in\bbQ[T]$ and $\deg X_k\leq k$; and $Y(\zeta X)=\zeta Y(X)$.
\end{lem}

\begin{proof}
Write $Y=X\,y(T)$, $T=X^5$, $y=1+\sum_{k\geq1}u^kY_k(T)$. Then $V/Y+Y^4/4=X^{-1}\bigl(V/y+Ty^4/4\bigr)$, and by \cref{lem:con-form}~(c) the equation becomes $V/y+Ty^4/4=\sum_{m\geq1}\frac{u^{m-1}}m\bigl(1+\frac{T^m}{5m-1}\bigr)$. At $u^0$ it holds with $V=1$, $y=1$. At $u^k$, $k\geq1$, it reads $(T-1)Y_k+V_k=R_k(T)$, where $V_k:=[u^k]V$ and $R_k$ is a polynomial built from $Y_1,\dots,Y_{k-1}$, $V_1,\dots,V_{k-1}$ and the term $\frac1{k+1}(1+\frac{T^{k+1}}{5k+4})$. It has a solution with $Y_k$ a polynomial if and only if $V_k=R_k(1)$, and then $Y_k=(R_k(T)-R_k(1))/(T-1)$ is unique. By induction $\deg Y_j\leq j$ for $j<k$; the $u^k$-parts of $V/y$ and of $Ty^4/4$ other than $V_k-Y_k$ and $TY_k$ then have degree at most $k$ and $k+1$, so $\deg R_k\leq k+1$ and $\deg Y_k\leq k$. At $k=1$: $R_1=\frac12+\frac{T^2}{18}$, $V_1=\frac59$, $Y_1=(1+T)/18$. Since $\hat\cW'(1)=0$ and $Y'(1)\in1+u\bbQ\fps{u}$ is a unit, $W_V'(Y(1))=0$, i.e. $Y(1)^5=V$ with $Y(1)\in1+u\bbQ\fps{u}$, so $Y(1)=V^{1/5}$; then $W_V(Y_c)=\frac54V^{4/5}$, and $\hat\cW(1)=\frac54\tau^{-1}$ gives $V=\tau^{-5/4}$. In the variables $x=sX$ and $sY$ the map reads $sY=x\bigl(1+\sum c_{k,d}u^{k-d}x^{5d}\bigr)$ with $d\leq k$, a change of variables in $\bbQ\fps{u,x}$ of the form $x\,(1+g(u,x^5))$ with $g(0,0)=0$. Its compositional inverse exists in $\bbQ\fps{u,y}$ and commutes with $y\mapsto\zeta y$, so it has the form $y\,(1+h(u,y^5))$ with $h(0,0)=0$; a monomial $u^ay^{5d}$ of $h$ is $u^{a+d}Y^{5d}$ at $y=sY$, and returning to $X$, $Y$ gives the stated form of the inverse with $\deg X_k\leq k$. Finally $T=X^5$ is invariant under $X\mapsto\zeta X$.
\end{proof}

\begin{lem}[Reduction modulo exact forms]\genusstatus{AI-proposed, unchecked}
\label{lem:con-red}
Let $\fa_Y(Y;\zh):=\hat\fa(X(Y);\zh)\,X'(Y)$, the hatted amplitude in the variable $Y$ of \cref{lem:con-cfu}; it lies in $\bbQ[Y,\zh]\fps{u}$. For $n\geq4$,
\[
\zh\,\dd\bigl(Y^{n-3}\re^{-W_V/\zh}\bigr)=\bigl((n-3)\zh Y^{n-4}+VY^{n-5}-Y^n\bigr)\re^{-W_V/\zh}\,\dd Y,
\]
and there are $\pi_m\in\bbQ[\zh]\fps{u}$ ($m=-1,\dots,3$) and $\mathcal E\in Y\bbQ[Y,\zh]\fps{u}$ with
\[
\fa_Y\,\re^{-W_V/\zh}\,\dd Y=\sum_{m=-1}^{3}\pi_m\,Y^m\re^{-W_V/\zh}\,\dd Y+\dd\bigl(\mathcal E\,\re^{-W_V/\zh}\bigr).
\]
With $\deg Y=1$ and $\deg\zh=-1$, $\pi_m$ contains only powers $\zh^k$ with $k\equiv m\bmod5$; for $m=\nu_b$ this is $k\equiv-b\bmod5$.
\end{lem}

\begin{proof}
The displayed identity is $\frac{\dd}{\dd Y}(Y^{n-3}\re^{-W_V/\zh})=\bigl((n-3)Y^{n-4}-Y^{n-3}W_V'(Y)/\zh\bigr)\re^{-W_V/\zh}$ with $W_V'=Y^3-VY^{-2}$; so $Y^n\equiv(n-3)\zh Y^{n-4}+VY^{n-5}$ modulo exact forms. At each $u$-order, $\fa_Y$ is a polynomial in $Y$ and $\zh$, because $\hat\fa\in\bbQ[X,\zh]\fps{u}$ and the inverse map of \cref{lem:con-cfu} is polynomial at each $u$-order. Each reduction step lowers the $Y$-degree by at least $4$, so finitely many steps bring a polynomial to the exponents $-1,\dots,3$, the lowest exponent produced being $4-5=-1$; the powers of $V$ bring power series in $u$, and each $u$-order receives finitely many steps. The primitives are $Y^{n-3}$ with $n\geq4$, whence $\mathcal E\in Y\bbQ[Y,\zh]\fps{u}$. Grading: $W_V(\zeta Y)=\zeta^{-1}W_V(Y)$; $\hat\fa(\zeta X;\zh)=\hat\fa(X;\zeta\zh)$, because $\fa$ depends on $x$ only through $xz$ and $x^5$; and the uniformising map commutes with $X\mapsto\zeta X$. So every monomial $Y^a\zh^b\,\dd Y$ of $\fa_Y\,\dd Y$ has $a-b+1\equiv1\bmod5$, and the relation preserves the degree modulo $5$. Hence every monomial $\zh^kY^m\,\dd Y$ of the result has $m-k\equiv0$; and $\nu_b\equiv-b\bmod5$ (\cref{def:set-model}). Uniqueness of the $\pi_m$ is not needed: \cref{lem:con-id} below determines them.
\end{proof}

% Section 8 (part 6): the model identity; signs, weights and identification; the relabelling.

\begin{lem}[The model identity]\genusstatus{AI-proposed, unchecked}
\label{lem:con-S2}
Let $W_1(\tilde p):=\tilde p^4/4+1/\tilde p$. For every $b\in\bbZ/5$, $\SPn_{\tilde p=1}[-W_1;\tilde p^{\nu_b}](t)=\varepsilon_b\Fs_b(t)$; equivalently, the full stationary phase $\mathrm{SP}_{\tilde p=1}[\tilde p^{\nu_b}\re^{-W_1(\tilde p)/t}\,\dd\tilde p]$ divided by $(2\pi t/5)^{1/2}\re^{-5/(4t)}$ is $\varepsilon_b\Fs_b(t)$.
\end{lem}

\begin{proof}
(The letters $J_m$, $G_b$, $U$ are local.) Let $A:=\bbQ\fps{v-1}$ and $\varsigma:=v^{1/5}\in A$. The phase $-\Wv(p)$ has the critical point $p=\varsigma$ (since $\Wv'(p)=p^3-v/p^2$) with $\Wv''(\varsigma)=5\varsigma^2$, a unit of $A$. Let $J_m:=\mathrm{SP}_{p=\varsigma}[p^m\re^{-\Wv(p)/z}\,\dd p]$, the full stationary phase over $A$. By \cref{lem:fsp-dq}, with the moving critical point $\varsigma$, $z\,\dd J_m/\dd v=-J_{m-1}$; by \cref{lem:fsp-ibp} applied to $\dd(\re^{-\Wv/z})$, $J_3=vJ_{-2}$. Hence $(J_0,-J_{-1},-J_3,J_2,J_1)=(\varepsilon_bJ_{\nu_b})_b$ solves the model system $z\,\dd Y_b/\dd v=b^{\mathrm m}_bY_{b+1}$, $b^{\mathrm m}=(1,-1/v,1,-1,-1)$ (\cref{prop:set-modrows}); the row $b=1$ uses $J_{-2}=J_3/v$, the other rows the derivative identity alone. Next, \cref{lem:fsp-cov} over $A$ with the linear map $p=\varsigma\tilde p$, for which $\Wv(\varsigma\tilde p)=\varsigma^4W_1(\tilde p)$, followed by \cref{lem:fsp-homog} with the unit $\mu=\varsigma^4$, gives
\[
J_{\nu_b}=(2\pi z/5)^{1/2}\re^{U/z}\varsigma^{\nu_b-1}G_b(z/\varsigma^4),\qquad U=-\tfrac54\varsigma^4,
\]
where $G_b(t):=\SPn_{\tilde p=1}[-W_1;\tilde p^{\nu_b}](t)$ does not depend on $v$. With $\dd U/\dd v=-1/\varsigma$ and $\dd\varsigma/\dd v=1/(5\varsigma^4)$ the chain rule gives $z\frac{\dd}{\dd v}\bigl[\re^{U/z}\varsigma^{k}F(z/\varsigma^4)\bigr]=\re^{U/z}\varsigma^{k-1}\mathcal T_kF$ with $\mathcal T_k=-1+\frac t5(k-4t\frac{\dd}{\dd t})$, so the system becomes the colour-$0$ T-system of \cref{prop:set-modrows} for $F_b:=\varepsilon_bG_b$. Since $G_0(0)=1$, the uniqueness half of \cref{prop:set-modrows} (pivot $4n$) gives $F_b=\Fs_b$, i.e. $G_b=\varepsilon_b\Fs_b$.
\end{proof}

\begin{lem}[Signs, weights, identification]\genusstatus{AI-proposed, unchecked}
\label{lem:con-id}
Work in a conifold frame of \cref{def:set-frames}, and let $Y_c=\tau^{-1/4}$ and the $\pi_m$ be as in \cref{lem:con-cfu,lem:con-red}.
\begin{enumerate}
\item[(a)] $\mathrm{SP}_{Y_c}[Y^{\nu_b}\re^{-W_V/\zh}\,\dd Y]=N_{\zh}\,\varepsilon_b\,\tau^{(1-\nu_b)/4}\Fs_b(\tau\zh)$, where $N_{\zh}:=(2\pi\zh/5)^{1/2}\re^{-W_V(Y_c)/\zh}$ and $W_V(Y_c)=\hat\cW(1)$.
\item[(b)] $\PhiSP=\sum_b\varepsilon_b\,\pi_{\nu_b}\,\tau^{(1-\nu_b)/4}\Fs_b(\tau\zh)$.
\item[(c)] Let $P$ be a graded regular gauge with coefficients in $\bbQ\fps u$ from the pulled-back model to the quintic system (\cref{def:gau-onepole,prop:gau-graded}), and let $\Phi_P$ be $\varpi_0$ times the colour-$0$ entry of the hatted row $0$ of $P\,\Qr^{\mathrm m}(v(u))$; for $P=\Pc$ this is $\Phican$ (\cref{thm:gau-R0}). Then
\[
\Phi_P/\varpi_0=\sum_b\tau^{(1-\nu_b)/4}\,G_b\,\Fs_b(\tau\zh),\qquad G_b:=s^{\nu_b}P_{0b}(u,s^4\zh),
\]
where $G_b\in\bbQ\fps{u}\fps{\zh}$ contains only powers $\zh^k$ with $k\equiv-b\bmod5$, and $[\zh^k]G_b$ is divisible by $u^{(\nu_b+4k)/5}$.
\item[(d)] Let $\Delta_b\in\bbQ\fps{u}\fps{\zh}$, $b\in\bbZ/5$, be series that contain only powers $\zh^k$ with $k\equiv-b\bmod5$. If $\sum_b\tau^{(1-\nu_b)/4}\Delta_b\Fs_b(\tau\zh)=0$, then every $\Delta_b$ vanishes.
\end{enumerate}
Hence, if $P$ is a graded regular gauge with $\Phi_P=\PhiSP$, then $G_b=\varepsilon_b\pi_{\nu_b}/\varpi_0$ for every $b$.
\end{lem}

\begin{proof}
(a) Put $Y=Y_c\tilde p$. Since $V=Y_c^5$, $W_V(Y_c\tilde p)=Y_c^4W_1(\tilde p)$, and $\zh/Y_c^4=\tau\zh$. By \cref{lem:fsp-cov} (a linear map over $\bbQ\fps{u}$) and \cref{lem:fsp-homog} with $\mu=Y_c^4=\tau^{-1}$, the left side is $Y_c^{\nu_b+1}\,\mathrm{SP}_{\tilde p=1}[\tilde p^{\nu_b}\re^{-W_1/t}\,\dd\tilde p]$ at $t=\tau\zh$, which by \cref{lem:con-S2} is $Y_c^{\nu_b+1}(2\pi\tau\zh/5)^{1/2}\re^{-5/(4\tau\zh)}\varepsilon_b\Fs_b(\tau\zh)$. Here $Y_c^{\nu_b+1}\tau^{1/2}=\tau^{(1-\nu_b)/4}$ and $\frac5{4\tau}=\frac54Y_c^4=W_V(Y_c)=\hat\cW(1)$ (\cref{lem:con-cfu}).

(b) By \cref{lem:fsp-cfu} the full stationary phase at $X=1$ of $\re^{-\hat\cW/\zh}\hat\fa\,\dd X$ equals that at $Y_c$ of its image $\re^{-W_V/\zh}\fa_Y\,\dd Y$. By \cref{lem:con-red} and \cref{lem:fsp-ibp}, applied once to the exact form $\dd(\mathcal E\re^{-W_V/\zh})$, whose primitive $\mathcal E$ is one amplitude over $\bbQ\fps{u}$, and by the $\bbQ\fps{u}\fps{\zh}$-linearity of the functional (\cref{def:fsp-wick}), this is $\sum_m\pi_m\,\mathrm{SP}_{Y_c}[Y^m\re^{-W_V/\zh}\,\dd Y]$. The exponents $m=-1,\dots,3$ are the $\nu_b$; insert (a) and divide by $N_{\zh}=(2\pi\zh/5)^{1/2}\re^{-\hat\cW(1)/\zh}$ (\cref{def:con-PhiSP}). The square roots agree: \cref{lem:fsp-cfu} takes $\hat\cW''(1)^{1/2}=Y'(1)\,W_V''(Y_c)^{1/2}$, (a) takes $W_V''(Y_c)^{1/2}=\sqrt5\,Y_c$, and \cref{def:con-PhiSP} takes $\hat\cW''(1)^{1/2}=\sqrt5\,(1-u)^{-1/2}$. Differentiating $\hat\cW=W_V\circ Y$ twice at $X=1$, where $W_V'(Y_c)=0$, gives $5/(1-u)=5Y_c^2\,Y'(1)^2$; as $Y_c$, $Y'(1)$ and $(1-u)^{-1/2}$ lie in $1+u\bbQ\fps u$, where a square root with constant term $1$ is unique, $Y'(1)\,Y_c=(1-u)^{-1/2}$.

(c) The colour-$0$ model column is $\Qr^{\mathrm m}_b=\varsigma^{\nu_b-1}\Fs_b(z/\varsigma^4)$ after the common exponential is stripped (\cref{prop:set-modrows,thm:gau-R0}), and $\varsigma=s\tau^{-1/4}$, $z/\varsigma^4=\tau\zh$ (\cref{def:set-tau}). Row $0$ of $P\Qr^{\mathrm m}$, hatted by $s^{1-\nu_0}=s$, is $\sum_bs\,P_{0b}\,s^{\nu_b-1}\tau^{(1-\nu_b)/4}\Fs_b(\tau\zh)$, which is the stated formula. For graded $P=\sum_kz^k\sum_{b'}p_{k,b'}E_{b'+k,b'}$ the entry $P_{0b}$ contains only $z^k$ with $k\equiv-b$, and $s^{\nu_b}z^k=u^{(\nu_b+4k)/5}\zh^k$ with $\nu_b+4k\equiv-b-k\equiv0\bmod5$ and $\nu_b+4k\geq0$.

(d) If not all $\Delta_b$ vanish, let $N$ be the least $\zh$-degree occurring in them. Only $b\equiv-N$ can contribute to $\zh^N$, through $\tau^{(1-\nu_b)/4}\varepsilon_b[\zh^N]\Delta_b$, because $\Fs_b(0)=\varepsilon_b$; as $\tau$ is a unit and $\varepsilon_b=\pm1$, this coefficient vanishes, a contradiction.

Finally, if $\Phi_P=\PhiSP$, subtract (b) from $\varpi_0$ times (c): the family $\Delta_b:=G_b-\varepsilon_b\pi_{\nu_b}/\varpi_0$ is graded as in (d), by (c) and \cref{lem:con-red} ($\nu_b\equiv-b$), so it vanishes.
\end{proof}

\begin{lem}[The relabelling]\genusstatus{AI-proposed, unchecked}
\label{lem:con-relab}
Let $\hat\cV:=\Res_X[\re^{-\hat\cW/\zh}\hat\fa\,\dd X]$, computed coefficientwise in $u$ in the residue ring $\cR_{\mathrm{res}}$ of \cref{lem:fsp-res}; at each $u$-order the integrand is $\re^{-(1/X+X^4/4)/\zh}$ times a Laurent polynomial in $X$ over $\bbQ[\zh,\zh^{-1}]$, an element of $\cR_{\mathrm{res}}$. Then $\cV=s\,\hat\cV$ coefficientwise: $[u^N\zh^j]\hat\cV=[u^{N+(1-4j)/5}z^j]\cV$ for all $N$ and $j$.
\end{lem}

\begin{proof}
\emph{Labels.} Monomials are $u^Ez^\beta x^\alpha$ in the flat variables and $u^N\zh^jX^\alpha$ in the hatted ones; the substitution $x=sX$, $z=s^4\zh$, $u=s^5$ acts on exponents by $\iota(E,\beta,\alpha)=(E+(4\beta+\alpha)/5,\beta,\alpha)$, which is injective. With $\dd x=s\,\dd X$ the claim is the coefficientwise statement.
\emph{The family.} The logarithm of $\re^{-\cG/z}\fa$ is a sum of atoms $-x^{5m-1}/(m(5m-1)z)$, $x^{5n}/(2n)$ and $\bk_\kappa\,c_{\kappa,P}\,(xz)^{2\kappa-1}x^{5P}$ ($\kappa\geq1$ a local index), with $c_{\kappa,P}:=5^{1-2\kappa}\delta_{P,0}-5\binom{2\kappa-2+5P}{5P}$ (\cref{lem:con-form}~(b)). A member of the family is a triple ($k$, $r$, a finite multiset of atoms); its term is $[u^r](\Xi/u)^k\,\frac{(-1)^k}{k!}\,u^{k+r}(xz)^{-k}$ times the product of the atoms divided by the factorials of the multiplicities. By \cref{def:con-PhiSP}, $\cV$ is the sum of the $x^{-1}$-parts of these terms.
\emph{Flat fibres.} A fixed $(u^E,z^\beta,x^{-1})$ receives finitely many members: $k\leq E$, and every atom has positive $x$-degree, so the multiset is bounded by the $x$-degree $k-1$.
\emph{Hatted fibres.} Every atom $x^az^b$ has hatted $u$-order $(a+4b)/5\geq0$, the only atom of hatted order $0$ is $-x^4/(4z)\mapsto-X^4/(4\zh)$, and $u^{k+r}(xz)^{-k}\mapsto u^r(X\zh)^{-k}$. Fix $(u^N,\zh^j,X^{-1})$. Then $r\leq N$, at most $N$ atoms have positive hatted order, from a finite list, and the $X$- and $\zh$-balances $-k+4l+a'=-1$, $-k-l+b'=j$, where $l$ is the multiplicity of $X^4/(4\zh)$ and $a'$, $b'$ are the degrees of the other atoms, determine $k$ and $l$ (the determinant is $5$). So this fibre is finite too.
\emph{$\hat\cV$ is the hatted sum.} In $\cR_{\mathrm{res}}\fps{u}$ the product is the convergent sum of the products of coefficients, and $\re^{-(\Xi/u)/(X\zh)}=\re^{-1/(X\zh)}\re^{-(\Xi/u-1)/(X\zh)}$ holds coefficient by coefficient: at $(u^N,X^{-n})$ both sides equal $(-1/\zh)^n[u^N](\Xi/u)^n/n!$, a finite expression. So $\hat\cV$ is the sum of the family in the hatted order.
\emph{Conclusion.} $\cV$ and $s\hat\cV$ are sums of one family over the fibres of two injective labellings related by $\iota$, with finite fibres in both orders; so they agree coefficient by coefficient.
\end{proof}

% Section 8 (part 7): the key lemma, the residue, Theorem P(b).

\begin{lem}[The key lemma]\genusstatus{AI-proposed, unchecked}
\label{lem:con-key}
Expand $\Sigma:=\sum_b\pi_{\nu_b}\,s^{-\nu_b}\,\IC_{\nu_b}(v(u),z)$ into monomials $u^Ez^\beta$, writing $\zh=s^{-4}z$ in $\pi_m$,
\[
\IC_m(v,z)=\sum_{i'\geq0}\frac{(-1/(4z))^{i'}}{i'!}\,\frac{(-v/z)^e}{e!},\qquad e:=m+4i'+1\quad(\text{terms with }e\geq0),
\]
and $v^e=u^eV^e$ (\cref{def:gau-van}; $v(u)=u\,\tau^{-5/4}$ by \cref{def:set-tau}, and $V=\tau^{-5/4}$ by \cref{lem:con-cfu}). A monomial $u^i\zh^j$ of $\pi_m$, the $i'$-th term of $\IC_m$ and a monomial $u^l$ of $V^e$ contribute to $u^Ez^\beta$ with $\beta=j-5i'-m-1$ and $E=i+l-(4\beta-1)/5$. Hence only $\beta\equiv4\bmod5$ occurs, and the coefficient of $u^Ez^\beta$ receives contributions only from the $u^i$-coefficients of the $\pi_m$ with $i\leq E+(4\beta-1)/5$, through finitely many terms; $\Sigma$ is a well-defined formal sum of monomials $u^Ez^\beta$.
\end{lem}

\begin{proof}
The $z$-exponent is $j-i'-e=j-5i'-m-1$, and $j\equiv m\bmod5$ by \cref{lem:con-red}, so $\beta\equiv-1$. The $u$-exponent is $i-\frac{4j}5-\frac m5+e+l$; substituting $j=\beta+5i'+m+1$ and $e=m+4i'+1$ gives $E=i+l-(4\beta-1)/5$. So $i=E+(4\beta-1)/5-l\leq E+(4\beta-1)/5$; for fixed $(E,\beta)$ and $i$ the exponent $l$ is fixed, $\pi_m$ is a polynomial in $\zh$ at the order $u^i$, so finitely many $j$ occur, and then $i'$ is fixed.
\end{proof}

\begin{prop}[The residue]\genusstatus{AI-proposed, unchecked}
\label{prop:con-res}
Work in a conifold frame, and let $P$ be a graded regular gauge with coefficients in $\bbQ\fps u$ with $\Phi_P=\PhiSP$ (\cref{lem:con-id}). Then
\[
\cV=\varpi_0\,\bigl(P\,\YC(v(u),z)\bigr)_0\qquad\text{in }\bbQ\lser{z}\fps{u}.
\]
\end{prop}

\begin{proof}
By \cref{lem:con-relab}, $\cV=s\hat\cV$ coefficientwise. By \cref{lem:fsp-res} the residue on $\cR_{\mathrm{res}}\fps{u}$ is invariant under the uniformising map, termwise in $u$, and vanishes on exact forms; with \cref{lem:con-red},
\[
\hat\cV=\Res_Y\bigl[\re^{-W_V/\zh}\fa_Y\,\dd Y\bigr]=\sum_{m=-1}^3\pi_m\,\Res_Y\bigl[Y^m\re^{-W_V/\zh}\,\dd Y\bigr],
\]
the products being formed in the hatted variables, with finite fibres, because $\pi_m$ is a polynomial in $\zh$ at each $u$-order. The model residues are, term by term,
\[
\Res_{Y=0}\bigl[Y^m\re^{-W_V/\zh}\,\dd Y\bigr]=\sum_{i'\geq0}\frac{(-1/(4\zh))^{i'}}{i'!}\,\frac{(-V/\zh)^e}{e!}=s^{-(m+1)}\IC_m(v(u),z),\qquad e=m+4i'+1:
\]
the substitution $p=sY$, with $v=uV$ and $z=s^4\zh$, turns $\Wv(p)/z$ into $W_V(Y)/\zh$, and expanding $V^e$ in $u$ is a relabelling of the kind treated in \cref{lem:con-relab}. The exponent map $\iota$ of \cref{lem:con-relab} is additive, so the image of a product with finite fibres is the product of the images. Hence $\cV=s\sum_m\pi_m\,s^{-(m+1)}\IC_m=\Sigma$, the formal sum of \cref{lem:con-key}, coefficient by coefficient. By \cref{lem:con-id}, $\pi_{\nu_b}=\varepsilon_b\varpi_0G_b=\varepsilon_b\varpi_0s^{\nu_b}P_{0b}$, and $\YC_b=\varepsilon_b\IC_{\nu_b}$ (\cref{def:gau-van}); the signs cancel, $\varepsilon_b^2=1$, and so do the powers of $s$, $s\cdot s^{-(\nu_b+1)}\cdot s^{\nu_b}=1$; so $\Sigma=\varpi_0\sum_bP_{0b}\YC_b=\varpi_0(P\YC)_0$. No power of $\tau$ survives: the residue side carries none, and $G_b$ is defined without it. Since $\cV\in z^{-1}\bbQ[z^{-1}]\fps{u}$ (\cref{def:con-PhiSP}), the identity holds in $\bbQ\lser{z}\fps{u}$.
\end{proof}

\begin{thm}[Theorem P(b)]\genusstatus{AI-proposed, unchecked}
\label{thm:con-Pb}
In every conifold frame of \cref{def:set-frames}, $\PhiSP=\Phican$ in $\bbQ\fps{u}\fps{\zh}$.
\end{thm}

Besides the lemmas cited in its steps, the proof uses Theorem~U (\cref{thm:gau-U}) for the pair of systems of \cref{prop:gau-quintic,prop:gau-equalprod}, whose hypotheses hold in the conifold frames, the gauge identity \cref{thm:gau-G}, the normalisations $\varpi_0(0)=1$ and $\tc'(0)=1$ of the conifold frames, and $\mathrm{ph}^{\mathrm{can}}_0=1$ (\cref{lem:gau-diag}, through \cref{thm:gau-R0}~(ii)). It does not use (QRR).

\begin{proof}
Steps 1--4 are \cref{lem:con-cfu} (the uniformising map), \cref{lem:con-red} (the reduction modulo exact forms), the invariance of both functionals under the uniformising map and their vanishing on exact forms (\cref{lem:fsp-cfu,lem:fsp-ibp,lem:fsp-res}), and the identification of \cref{lem:con-S2,lem:con-id}.

\emph{Step 5 (regularity).} By \cref{thm:con-Pa}~(iii), $\PhiSP$ solves HPF; by \cref{thm:gau-R0}~(ii), which rests on \cref{cor:gau-cycsys}, so does $\Phican$; both have constant term $1$ and coefficients in $\bbQ\fps{u}$. By \cref{lem:gau-poly}, $\PhiSP=\Phican\,\mathsf s(z)$ for a series $\mathsf s\in1+z^5\bbQ\fps{z^5}$ (the letter $\mathsf s$ is local to this proof; $z^5=u^4\zh^5$). Put $P^{\mathrm{SP}}:=\mathsf s(z)\Pc$. By \cref{thm:gau-U}~(b) and \cref{prop:gau-graded} it is a graded regular gauge, and $\Phi_{P^{\mathrm{SP}}}=\mathsf s(z)\Phican=\PhiSP$. The gauge identity for the canonical gauge, $\Pc\YC(v(u),z)=\YCQ$ (\cref{thm:gau-G}), multiplied by $\mathsf s(z)$, gives $(P^{\mathrm{SP}}\YC)_0=\mathsf s(z)(\YCQ)_0$.

\emph{Step 6 (the residue).} By \cref{prop:con-res} for $P=P^{\mathrm{SP}}$, $\cV=\varpi_0\,\mathsf s(z)\,(\YCQ)_0$ in $\bbQ\lser{z}\fps{u}$.

\emph{Step 7 (the order $u^1$).} On the left only $k=1$ contributes to $u^1$ in \cref{def:con-PhiSP}, because $\Xi=u+O(u^2)$, and $[x^0](\re^{-\cG/z}\fa)=1$, because $\cG=O(x^4)$ and $\fa(0,z)=1$; so $[u^1]\cV=-1/z$. On the right, $[u^0](\YCQ)_0=0$, since $y^{(0)}=-e_1$ and the terms $y^{(n)}$, $n\geq1$, vanish at $u=0$; and $[u^1](\YCQ)_0=-1/z$, since by \cref{lem:gau-vanq} only $y^{(1)}_0=-\tc$ contributes at $u^1$ and $\tc=u+O(u^2)$. With $\varpi_0(0)=1$ the right side has $u^1$-coefficient $-\mathsf s(z)/z$. Hence $\mathsf s=1$, and $\PhiSP=\Phican$.
\end{proof}

% Section 8 (part 8): the glue. Last part of Section 8.

\subsection{The glue}
\label{sec:con-glue}

For a colour-$0$ WKB function with polynomial coefficients write $\Phi(z,1):=\sum_n\mathrm{ph}_n(1)z^n$, its value at $u=1$, where $\zh=z$.

\begin{thm}[The glue]\genusstatus{AI-proposed, unchecked}
\label{thm:con-glue}
In $\bbQ(\zeta)\fps{z}$,
\begin{equation}
\label{eq:con-glue}
h(z)\,\Dn_0(z)=\Phican(z,1);
\end{equation}
hence $\Dn_0\in1+z\bbQ\fps{z}$, and, by \cref{thm:con-Pb}, $\Dn_0(z)=\PhiSP(z,1)/h(z)$.
\end{thm}

\begin{proof}
(a) By \cref{thm:gau-R0}~(ii)--(iv), the canonical flat amplitude $A^{\mathrm{can}}=\sum_nA^{\mathrm{can}}_nz^n$ has $A^{\mathrm{can}}_n=L^{4n+1}\mathrm{ph}^{\mathrm{can}}_n(L^{-5})\in\bbQ[L,L^{-1}]$, the $\mathrm{ph}^{\mathrm{can}}_n$ being polynomials. In $\bbQ(\zeta)\fps{q}$, with $L=(1-5^5q)^{-1/5}$, $A^{\mathrm{can}}$ and the actual amplitude $A^{\mathrm{act}}$ are colour-$0$ WKB columns of $\PFfive$ with the same exponential $\re^{u_0/z}$ and the same leading term $L$, so by formal WKB uniqueness (\cref{lem:gau-wkb}, with $\Bbbk=\bbQ(\zeta)$) they differ by a constant series in $1+z\bbQ(\zeta)\fps{z}$; hence every $A^{\mathrm{act}}_n$ lies in $\bbQ(\zeta)[L,L^{-1}]$. The ring $\bbQ(\zeta)[L,L^{-1}]\fps{z}$ embeds in $\bbQ(\zeta)\lser{s}\fps{\zh}$, where the comparison of \cref{thm:gau-R0} reads $A^{\mathrm{act}}\,\Dn_0=A^{\mathrm{can}}$; so the constant series is $\Dn_0$. (b) Evaluation at $L=1$, that is $u=1$, $q=0$, $\zh=z$, is a ring homomorphism $\bbQ(\zeta)[L,L^{-1}]\to\bbQ(\zeta)$; applied coefficientwise it gives $A^{\mathrm{act}}(z)|_{L=1}\,\Dn_0(z)=\Phican(z,1)$. As $(1-5^5q)^{-1/5}=1+O(q)$, $A^{\mathrm{act}}_n|_{L=1}$ is the coefficient of $q^0$ of $A^{\mathrm{act}}_n$, which is $h_n$ by (QRR) in the form of \cref{thm:set-qrr}, \eqref{eq:set-qrr}. As $h\in1+z\bbQ\fps{z}$ is invertible, $\Dn_0=\Phican(z,1)/h(z)$ has rational coefficients, and \cref{thm:con-Pb} gives the second form. The inputs used here are the value at $q=0$ of (QRR) and the actual rows with their WKB form, which rest on (M), (G0) (\cref{thm:set-mirror,thm:set-G0}) and the flatness part of (QRR) (\cref{def:set-rows}~(b)) through \cref{cor:set-cyclicPF5} and \cref{thm:gau-R0}; no other step of this section uses an input.
\end{proof}

\begin{cor}\genusstatus{AI-proposed, unchecked}
\label{cor:con-A2}
$\Dn=1$ if and only if $\PhiSP(z,1)=h(z)$. By \cref{thm:con-Pa}~(iv), $\PhiSP(z,1)=c(z;1)$, the value at $L=1$ of the Laurent polynomials $c_n$; so the condition reads $c(z;1)=h(z)$ (\cref{cor:gam-equiv}).
\end{cor}

\begin{proof}
By \cref{cor:gau-colours}, $\Dn_i(z)=\Dn_0(\zeta^{-i}z)$, so $\Dn=1$ if and only if $\Dn_0=1$; by \eqref{eq:con-glue} and \cref{thm:con-Pb}, $\Dn_0=1$ if and only if $\PhiSP(z,1)=h(z)$. By \cref{thm:con-Pa}~(iv), $c_n(1)=\mathrm{ph}^{\mathrm{SP}}_n(1)$.
\end{proof}



